\documentclass{article}
\usepackage{amsmath}
\begin{document}
\title{Two camera angles can restore a covered cue, but only if they fit the capture deadline}
\maketitle

\paragraph{Changed detector law.}
Retain the earlier authored far scene and initially reach commanded
lateral offset $.2$ at rest by time $.2$. Actual offset has fixed
unknown bias $\pm .015$. The camera's first angle is installed
in advance, with a movable yaw actuator distinct from the formation
control. During each useful exposure both pose and angle are fixed.
The receiver must positively identify a far return through a named
ray in each applicable phase; it may not infer a world from failure
to see the target. The scene, target, identity and transport binding
are fixed during these two exposures. Each exposure may take $.125$
units and the modeled extra settling allowance is $.01$ total. This
is an explicit alternative detector premise: two exposures with
the original full $.27794$ dwell cannot fit. No measured integration
time, station yaw actuator or optics are supplied by these premises.

\paragraph{Exact geometric and slew limits.}
For the two target depths $37,38$ and panel at 8, a small positive
world ray slope $s$ has the open common positive aperture
\[
 (.20-8s,\ .28-38s).
\]
For arbitrary rays the same inequalities are necessary for a
positive-side ray hitting both worlds; other placement constraints
could only shrink the interval. An actual uncertainty interval
cannot cross the central region where both worlds have no visible
target. For nominal $b\le .2$ and symmetric radius $e$, one ray
must cover both $b-e$ and $b+e$. This needs
$s\ge(.2-b+e)/8$ and $s\le(.28-b-e)/38$ with strictness at the
boundaries. The maximum possible radius supremum is attained for
nominal $.2$ in the limiting sense and equals $0.64/46$, about
13.913 mm. The previous fixed lattice remains more restrictive.
With two rays sorted by slope, the high slope must reach the
left endpoint and the low slope the right endpoint. The intervals
need overlap, so $38s_{\rm high}-8s_{\rm low}\le .08$. Optimizing
at nominal $.2$ (decreasing $b$ worsens the combined inequality)
gives the no-slew radius supremum
\[
 e\le .08(46)(8)/(38^2+8^2)\simeq19.523\text{ mm}.
\]
At its endpoint the two intervals touch; any smaller nearby
radius has a strict constructive cover. These are restrictions on
positive single-ray evidence at each phase, not general bounds
on arbitrary full-frame inference or concurrent separate detectors.

At $b=.2,e=.015$, minimum angle change has the sharp infimum
\[
 J_*=\arctan(e/8)-\arctan((.08-e)/38)>.0001644\ {\rm rad}.
\]
Both endpoint constraints bound the respective slopes, and arctan
is monotone. The limiting apertures have strict overlap, so one
can approach this travel from above without losing coverage.
With acceleration 4 and speed limit $.01$ rad per unit,
a rest-to-rest slew needs at least $J_*/.01+.01/4>.01894$.
The available slate after exposure and settling is only
$q-.46\simeq .017943$ for $q=\sigma-.15$. Thus no such two-ray
schedule can guarantee the required capture under the slower
yaw limit at this nominal position.

\paragraph{Constructive timing arm.}
Use the slopes $.00190$ and $.00169$. They cover the whole
bias range with overlap, even with a separate fixed angular
calibration bias of magnitude $5\cdot 10^{-6}$ radians. Over
this interval the angle error changes slope by less than
$5.001\cdot 10^{-6}$, and the three strict endpoint/overlap
checks remain valid. This angular bias is fixed across the
rest-to-rest move; two unrelated angle errors would require
a separate travel and calibration account. At speed limit $.02$
and acceleration 4, the chosen angle difference is below $.00021$
and takes at most $.00021/.02+.02/4=.0155$.
The two exposures and guard therefore finish by $.4755<q$.
The original diagnostic slopes $.002$ and $.0015$ also cover
geometrically, but their slew misses this charged deadline even
without a speed bound (and takes nearly $.03$ at the faster
pinned speed limit). A direct controller informed of the same
apertures, scheduler and actuator envelope yields the same result.

\paragraph{Source execution and remaining obligations.}
Run actual yaw projections before the two pinned native shader
passes, with modeled poses and independent held still angles.
The 168 sampled frames cross both schedules, seven lateral states,
three angle-bias values and two true worlds. The selected native
depth publication and scoped raw recognition are driven by the
rendered values. Select a qualifying as-of return from the first
completed eligible phase; if the second phase misses its cut, do
not credit it. Keep the same old depth source service, actual
old native radio callback, one-byte cue and exact receipt cutoff.
All 504 joined cases replay under calibrated logic; the stronger
arm has 126 timely cues, while the two failing arms still complete
some instances without guaranteeing the uncertainty envelope.
A native rendered image is not physical noise/visibility data,
actuator timing or an independently selected scene. No R1/R6
completion or architecture acceptance is claimed.
\end{document}
